\chapter{Covariates in the excitation: a Volterra system}
\label{ch:excitation}

\begin{motivation}
Sometimes the \emph{contagiousness} itself varies --- a disease spreads faster in
winter, a post is more re-shareable during a trend. Then covariates must modulate
the \emph{kernel}, not just the baseline. This is the genuinely hard case, and the
reason the whole Volterra viewpoint was worth developing: the governing equation
stops being a convolution. We show it is still solvable, exactly how, and at what
cost.
\end{motivation}

\section{The model}

\begin{model}[covariate-modulated excitation]\label{mod:ch9}
Let covariates $Z(t)\in\R^q$ scale the triggering strength,
\begin{equation}
\alpha_{}(t)=\kappa\theta\,\exp\!\big(\delta^\top Z(t)\big),\qquad
\phi(t,u)=\alpha(t)\,e^{-\theta(t-u)},
\end{equation}
so the time-varying branching ratio is $\kappa(t)=\kappa\,e^{\delta^\top Z(t)}$.
\end{model}

\section{The equation is no longer a convolution}

Averaging as before gives
\begin{equation}
\xi(t)=s(t)+\int_0^t K(t,u)\,\xi(u)\dd u,\qquad
K(t,u)=\kappa\theta\,e^{\delta^\top Z(t)}\,e^{-\theta(t-u)},
\label{eq:ch9-volterra}
\end{equation}
a \emph{general} (non-convolution) Volterra equation: $K(t,u)$ depends on $t$ and
$u$ \emph{separately}, not only on $t-u$.

\begin{whybox}[: what we lose, and why]
A convolution kernel means the system reacts to a stimulus the same way whenever it
arrives (time-invariance), which is exactly what let Laplace/Fourier diagonalise
it and gave a single transfer function (Ch.~\ref{ch:solving}, Ch.~\ref{ch:solvers}).
Once the gain rides on $Z(t)$, \emph{when} matters: no convolution theorem, no
constant transfer function, and the spectral/operator shortcuts no longer apply.
\end{whybox}

\begin{intuitionbox}
The same past event echoes more loudly in a high-$Z$ period than a low-$Z$ one. The
``filter'' is the same shape, but its volume knob is being turned over time.
\end{intuitionbox}

\section{What survives: existence, and a time-varying ODE}

\paragraph{Well-posedness survives.} \eqref{eq:ch9-volterra} is a standard linear
Volterra equation; by the general theory (Appendix~\ref{app:volterra}, which never
used convolution) it has a unique solution on every finite horizon. Existence does
not need time-invariance.

\paragraph{Exponential shape $\Rightarrow$ a finite LTV ODE.} The kernel is
\emph{separable}: $K(t,u)=\underbrace{\kappa\theta e^{\delta^\top Z(t)}e^{-\theta t}}_{a(t)}\,\underbrace{e^{\theta u}}_{u\text{ only}}$.
A separable Volterra kernel collapses to finitely many ODEs
(Appendix~\ref{app:volterra}). With the filtered history
$y(t)=\int_0^t e^{-\theta(t-u)}\xi(u)\dd u$,
\begin{equation}
y'(t)=\xi(t)-\theta\,y(t),\qquad
\xi(t)=s(t)+\kappa\theta\,e^{\delta^\top Z(t)}\,y(t),
\label{eq:ch9-ltv}
\end{equation}
a \emph{linear time-varying} (LTV) ODE: the same state as the convolution case, but
with a gain that breathes with the covariate.

\begin{intuitionbox}
``Same state, time-varying gain.'' We keep a finite-dimensional, integrable system;
we merely lose the constant-coefficient (hence transform-solvable) structure.
\end{intuitionbox}

\section{In code, and a fast exact path}

The general solver is the LTV integrator \code{solve\_mbpp\_ltv}; for a
piecewise-constant covariate the package uses an \emph{exact per-regime}
propagation (the gain is constant on each regime, so the 1-D linear ODE is solved
in closed form) --- fast and more accurate than a grid.

\begin{lstlisting}[caption={Forward solve and fitting with excitation covariates.}]
from hawkes_calibration import solve_mbpp_ltv, fit_mbpp_ic_excitation
# forward: xi for covariate-modulated excitation
xi = solve_mbpp_ltv(forcing, A0=[[kappa*theta]], B=[[theta]], t_grid,
                    modulation=lambda t: [[np.exp(delta @ Z(t))]])
# fit delta (the excitation-covariate effect) from interval counts:
res = fit_mbpp_ic_excitation(obs_times, counts, Z)   # recovers kappa, theta, mu, delta
\end{lstlisting}

\begin{examplebox}[: recovering $\delta$]
With $Z\equiv0$ the LTV solver reproduces the constant-kernel solver to $10^{-16}$;
a covariate that doubles excitation over a window elevates the intensity then
relaxes back. On MBPP-Poisson data the fitter recovers $\delta=1.2$ as $\approx1.14$
in $\sim$0.5\,s (the exact per-regime path); on genuine over-dispersed Hawkes data
(\code{exp10}) it is recovered on average ($1.49\pm0.9$ over groups).
\end{examplebox}

\begin{pitfall}
Excitation covariates make the IC-LL \emph{non-convex} (the $\delta$ enter through
$e^{\delta^\top Z}$ inside the intensity), and $\delta$ is weakly identified from
little data. Use many sequences and restarts; identification needs the covariate to
vary \emph{while events occur}. Multivariate excitation covariates add up to
$pM^2$ parameters and are harder still --- the supported fitter is univariate.
\end{pitfall}

\section{The general shape: Volterra quadrature}

If the triggering shape is not a finite sum of exponentials (e.g.\ power-law), the
kernel is not separable and no finite ODE exists. The package then solves
\eqref{eq:ch9-volterra} directly by trapezoid-Nyström quadrature
(\code{solve\_mbpp\_volterra}, Appendix~\ref{app:volterra}) --- $O(N^2)$ and
first-order, but fully general.

\begin{recap}
Excitation covariates turn the MBPP equation into a general Volterra equation:
well-posed, but no longer time-invariant, so the transform shortcuts vanish. For an
exponential shape it reduces to a finite LTV ODE (\code{solve\_mbpp\_ltv}, fast
exact per-regime path); for a general shape, to Volterra quadrature. Fitting
recovers $\delta$ but is non-convex and data-hungry. This is the precise content of
``excitation covariates are harder.''
\end{recap}

\exercise{Show that placing the modulation at the \emph{parent} time, $K(t,u)=
\kappa\theta e^{\delta^\top Z(u)}e^{-\theta(t-u)}$, is also separable, and write the
corresponding LTV ODE.}
\exercise{Why does the spectral (Fourier) operator of Ch.~\ref{ch:solvers} fail to
represent \eqref{eq:ch9-volterra} exactly, whereas it was exact for the
convolution case?}
